\documentclass[aps,pre,reprint,superscriptaddress,longbibliography,nofootinbib]{revtex4-2}

\usepackage{graphicx}
\usepackage{amsmath,amssymb,amsfonts}
\usepackage[utf8]{inputenc}
\usepackage{xcolor}
\usepackage{hyperref}
\usepackage{upgreek}
\usepackage{soul}
\usepackage[capitalise]{cleveref}


\usepackage[inline]{showlabels}
\renewcommand{\showlabelfont}{\small\ttfamily\color{lightgray}} % Custom label appearance

% Reference names. Every cross-reference in this file goes through \cref,
% so that section/equation/figure/appendix prefixes are generated rather
% than typed by hand.
\crefname{equation}{Eq.}{Eqs.}
\Crefname{equation}{Equation}{Equations}
\crefname{section}{Sec.}{Secs.}
\Crefname{section}{Section}{Sections}
\crefname{appendix}{Appendix}{Appendices}
\Crefname{appendix}{Appendix}{Appendices}
\crefname{figure}{Fig.}{Figs.}
\Crefname{figure}{Figure}{Figures}
\crefname{table}{Table}{Tables}
\Crefname{table}{Table}{Tables}

\hypersetup{
    colorlinks=true,
    linkcolor=[rgb]{0 0 0.5},
    urlcolor=[rgb]{0 0 0.5},
    citecolor= [rgb]{0 0 0.5}
}


\graphicspath{{figures/}}

\newcommand{\la}{\left\langle}
\newcommand{\ra}{\right\rangle}
\newcommand{\bs}{\boldsymbol}
\newcommand{\E}{\mathbf{E}}
\newcommand{\Var}{\text{Var}}
\renewcommand{\d}{\mathrm{d}}
\renewcommand{\u}{\bs{u}}
\newcommand{\q}{\bs{q}}
\newcommand{\Q}{\mathbf{Q}}
\newcommand{\R}{\mathbb{R}}
\newcommand{\n}{\bs{n}}
\newcommand{\e}{\bs{e}}
\newcommand{\T}[0]{\mathrm{T}}
\newcommand{\pd}[2]{\frac{\partial #1}{\partial #2}}
\newcommand{\tpd}[2]{\tfrac{\partial #1}{\partial #2}}


\newcommand{\ab}{{\alpha\beta}}
\newcommand{\mn}{{\mu\nu}}
\newcommand{\abmn}{{\alpha\beta\mu\nu}}

\newcommand{\eye}{\bs{\mathrm{I}}}
\renewcommand{\v}{\bs{v}}
\newcommand{\x}{\mathbf{x}}
\newcommand{\f}{\mathbf{f}}
\newcommand{\s}{\mathbf{s}}
\newcommand{\g}{\mathbf{g}}
\newcommand{\G}{\mathbf{G}}
\renewcommand{\H}{\mathbf{H}}
\newcommand{\F}{\mathbf{F}}
\newcommand{\I}{\mathbf{I}}
\newcommand{\J}{\mathbf{J}}
\newcommand{\bnabla}{\boldsymbol{\nabla}}
\newcommand{\A}{\mathbf{A}}
\newcommand{\B}{\mathbf{B}}
\newcommand{\U}{\mathbf{U}}
\newcommand{\bxi}{\mathbf{\xi}}
\newcommand{\p}{\mathbf{p}}
\newcommand{\W}{\mathbf{W}}
\newcommand{\bT}{\mathbf{T}}
\newcommand{\nh}{\hat{\bs{n}}}
\newcommand{\bdelta}{\bs\delta}
\renewcommand{\ss}{\mathrm{ss}}

\newcommand{\ee}{\mathrm{e}}
\newcommand{\ii}{\mathrm{i}}

\renewcommand{\k}{\bs{k}}
\renewcommand{\r}{\mathbf{r}}

\newcommand{\TT}{{\rm T}}
\newcommand{\Tr}{{\rm T}}


\newcommand{\lM}{l_{\rm M}}
\newcommand{\lm}{l_{\rm m}}

\newcommand{\elln}{\ell_{\rm n}}
\newcommand{\ellc}{\ell_{\rm c}}
\newcommand{\elld}{\ell_{\rm d}}

\newcommand{\calN}{\mathcal{N}}
\newcommand{\calF}{\mathcal{F}}
\newcommand{\calG}{\mathcal{G}}
\newcommand{\calK}{\mathcal{K}}
\newcommand{\rmR}{\mathrm{R}}

\providecommand{\Fpair}{\calF_{\mathrm{pair}}}
\providecommand{\Fent}{\calF_{\mathrm{ent}}}
\providecommand{\Tr}{\operatorname{tr}}
\newcommand{\frev}{f_{\rm rev}}
\newcommand{\xir}{\xi_{\rm r}}
\newcommand{\ecrit}{\varepsilon_{\rm c}}
\newcommand{\sigmarho}{\sigma_\rho^2}
\newcommand{\kB}{k_\mathrm{B}}
% Correction annotations (visible in PDF)
\newcommand{\corr}[1]{\textcolor{red}{\textbf{[CORR:}\ #1\textbf{]}}}
\newcommand{\corrm}[1]{{\color{red}#1}}

\renewcommand{\div}{\text{div}}

\newcommand{\pade}{\text{pade}}

\newcommand{\Pa}{\text{Pa}}
\newcommand{\minute}{\text{min}}
\newcommand{\micron}{\upmu\mathrm{m}}

\newcommand{\ash}[1]{\textcolor{blue}{\textbf{[ash:}\ #1\textbf{]}}}

% Quick visual marker for skeleton material -- strip \skel and \todo
% (already in definitions.tex) once real content replaces the bullets.
\newcommand{\skel}[1]{\textcolor{gray!70!black}{\itshape #1}}

% prompt, address all the prompts in the text
\newcommand{\prompt}[1]{\textcolor{blue!90!black}{[ #1 ]}}

% placeholder, will be filled later
\newcommand{\citeTODO}[1]{\textcolor{red}{[\textit{cite} #1]}}

% placeholder, will be filled later
\newcommand{\placeholder}[1]{\textcolor{purple}{[\textit{placeholder}] #1 }}
% Caption-safe bullet: \caption{} is a "moving argument" (also read by
% hyperref for PDF bookmarks and by \listoffigures), and an itemize
% environment inside it breaks with "Incomplete \iffalse" -- confirmed
% by direct test against this revtex4-2+hyperref combination. Use this
% macro for every bulleted caption instead of \begin{itemize}.
\newcommand{\capitem}{\par\noindent$\bullet$\ }

\begin{document}

\title{Fluctuating active nematics account for layer formation in \textit{Myxococcus xanthus} colonies}

\author{Ashot Matevosyan}
\author{Ricard Alert}

\date{\today}

\begin{abstract}
In colonies of \textit{Myxococcus xanthus}, cells leave the monolayer and build a second
layer preferentially at $+1/2$ topological defects of the cell-orientation field
\cite{copenhagen2021}. The mean compression a defect sustains falls three to four orders
of magnitude short of the pressure needed to extrude a cell, so the events cannot be the
mean crossing a threshold; they are rare excursions of a fluctuating stress, and these
monolayers carry traction fluctuations larger than the mean active stress
\cite{han2025local}. An active nematic has no fluctuation-dissipation relation to fix the
amplitude of its noise, so we obtain that amplitude from the microscopic dynamics rather
than supplying it: we coarse-grain $N$ self-propelled rods that reverse at rate $\frev$
into a Dean equation \cite{dean1996}, and integrate out the orientation to reach
stochastic equations for the density and the nematic order parameter. The reversal rate
fixes the mean active stress and the noise amplitude together, with no freedom between
them. Around a $+1/2$ defect these equations give a pressure variance modulated by about
three percent over a few micrometers. Because the extrusion threshold lies nineteen
standard deviations into the tail, that three percent is amplified into a
two-orders-of-magnitude enhancement of the layering probability, peaking off the defect
core, which is what Fig.~3f of Ref.~\cite{copenhagen2021} measures.
\end{abstract}

\maketitle

% ============================================================
\section{Introduction}
\label{sec:intro}
% Pages: ~1.5. Source: new. structure.md \S1.
% ============================================================
A monolayer of \textit{Myxococcus xanthus} on a wet surface is a dry active nematic. The
cells are rods a few micrometers long, they glide along their own axis, and momentum is
not conserved in the plane because the substrate takes it up, so the physics is carried
by the orientation field and not by a flow. Like any two-dimensional nematic the
monolayer is threaded by $\pm1/2$ disclinations. When the colony can no longer grow in
the plane, cells leave it, and the second layer they form does not nucleate uniformly:
the events concentrate at $+1/2$ defects, and not at $-1/2$ defects or in the bulk
\cite{copenhagen2021}. \Cref{fig:phenomenon} shows the phenomenon and the model we build
from it.

A deterministic explanation suggests itself. A $+1/2$ defect carries a polar active force
density, so it compresses the monolayer on one side of its core, and layer formation
might simply be where that compression is largest. The numbers rule this out. Extruding a
cell means peeling it off the substrate and pushing it through the water film that wets
the colony, and the surface tension of that film alone puts the threshold pressure at
some $3\times10^4~\Pa\cdot\micron$, while the mean compression a defect sustains in a
monolayer this stiff is three to four orders of magnitude smaller. Layer formation is
therefore not the mean reaching a threshold. It is a rare excursion of a fluctuating
stress, and fluctuations are what these monolayers have in abundance: the traction
fluctuations measured in Ref.~\cite{han2025local} exceed the mean active stress.

A rare-event rate needs the statistics of the stress and not its mean, and this is where
a continuum active-nematic theory \cite{marchetti2013hydrodynamics} stops being enough.
In equilibrium the amplitude of the noise is not free, because the
fluctuation-dissipation theorem ties it to the dissipation already present in the
equations. An active nematic is driven, has no equilibrium to relax to, and offers no
such relation. Noise added at the continuum level therefore arrives with an amplitude
that has to be supplied from outside, and the rate is exponentially sensitive to exactly
that amplitude. We do not supply it. We start instead from $N$ self-propelled rods that
reverse their gliding direction at a rate $\frev$, as \textit{M.\ xanthus} does,
coarse-grain them into a Dean equation \cite{dean1996} for the joint density in position
and orientation, and integrate out the orientation.

What comes out is a pair of stochastic equations in which the activity enters twice under
a single coefficient. The reversal noise generates an active stress with coefficient
$\zeta=\xi_0v_0^2/(4\frev)$, so that the activity is set by the distance $v_0/\frev$ a
cell glides between reversals, and the same
$\zeta$ sets the variance of the noise that survives coarse-graining. Mean stress and
fluctuation amplitude are not independent quantities in this model: one measured
biological parameter controls both, and neither can be adjusted without the other.

Solved around a $+1/2$ defect, these equations give a density variance made of two
pieces. One is uniform, the shot noise of a conserved flux resolved at the size of a
cell. The other follows the nematic texture, decays away from the core, and modulates the
local variance by about three percent. Three percent is small, and the reason it is
nonetheless decisive is the height of the barrier it acts on. The extrusion threshold
lies about nineteen standard deviations out in the tail of the pressure distribution, so
the logarithm of the crossing probability carries a prefactor $z_\infty^2/2\simeq180$,
and a three percent change in the variance moves the rate by a factor of a few hundred.
That exponential lever is the mechanism: a barely perceptible feature of the nematic
texture becomes a large, spatially structured change in a rate.

\Cref{sec:model} derives the stochastic equations. \Cref{sec:defect} solves them at a
$+1/2$ defect, first for the steady state the defect drifts in and then for the
fluctuations around it. \Cref{sec:nucleation} converts those fluctuations into a layering
probability and compares it with Ref.~\cite{copenhagen2021}. \Cref{sec:discussion}
collects what is fitted, what is not, and what would falsify the mechanism.

\begin{figure}[tbp]
  \centering
  \fbox{\parbox[c][4cm][c]{0.85\linewidth}{\centering\textbf{[FIGURE 1: TO MAKE]}}}
  \caption{%
\capitem Phenomenon $\to$ model, left to right or top to bottom: (a) experimental
      layer-formation event(s) at a $+1/2$ defect, adapted from \cite{copenhagen2021}
      (permission needed); (b) cartoon of a reversing self-propelled rod, rate $\frev$;
      (c) the coarse-grained fields $\rho(\r), \Q(\r)$ with a $+1/2$ defect marked.
    \capitem Candidate panel: \texttt{figures/myxo\_life.png} \cite{munoz2016} for
      colony-scale context (vegetative swarm $\to$ starvation $\to$ fruiting body).
      \todo{decide if this belongs here, is cut, or moves to a graphical-abstract-only role}.
    \capitem This is the one main-text figure that does not exist yet (structure.md,
      ``Open items'').
}
  \label{fig:phenomenon}
\end{figure}

% ============================================================
\section{Fluctuating active nematics from microscopic dynamics}
\label{sec:model}
% structure.md restructure, 2026-09-21. dean.tex 109-540.
% ============================================================

\subsection{The microscopic model}
\label{sec:model-micro}
% dean.tex 109-173 ("Microscopic equations of motion").
We start from Langevin equations for $N$ overdamped, self-propelled rods
\cite{cugliandolo2015,bertin2013}, each with position $\r_i$ and orientation $\theta_i$:
\begin{subequations}
  \label{eq:model}
\begin{align}
    \label{eq:langevin-position}
    \frac{\d\r_i}{\d t} &= \frac{\rho_0}{\xi_0}\F_i + v_0\,\p(\theta_i) k_i,
    \\
    \label{eq:langevin-angle}
    \frac{\d\theta_i}{\d t} &= \frac{\rho_0}{\xir}\Gamma_i + \lambda_i,
\end{align}
\end{subequations}
with $\p(\theta)=(\cos\theta,\sin\theta)$ the unit vector along the rod axis. Here $\xi_0/\rho_0$ and $\xir/\rho_0$ are
translational and rotational friction coefficients ($\rho_0$ a reference density
carried along for later convenience), and $\lambda_i$ is Gaussian white
noise, $\la\lambda_i(t)\lambda_j(t')\ra=2\Lambda\delta_{ij}\delta(t-t')$, with $\Lambda$
the rotational diffusion coefficient. The rods
interact through a pairwise potential $V(\r-\r',\theta,\theta')$, so that
$V_\mathrm{tot}=\tfrac12\sum_{i,j=1}^N V(\r_i-\r_j,\theta_i,\theta_j)$. The force and
torque on rod $i$ are the values at $(\r_i,\theta_i)$ of the fields
\begin{align}
    \label{eq:force-torque}
    \F(\r,\theta) &= -\pd{}{\r}\sum_{j=1}^N V(\r-\r_j,\theta,\theta_j),
    \nonumber\\
    \Gamma(\r,\theta) &= -\pd{}{\theta}\sum_{j=1}^N V(\r-\r_j,\theta,\theta_j),
\end{align}
with $N$ the number of rods. The sums run over all $N$ rods, including $j=i$: a pair
potential between identical rods satisfies
$V(\r-\r',\theta,\theta')=V(\r'-\r,\theta',\theta)$, and at $\r=\r'$ rotational
invariance leaves it a function of $\theta-\theta'$ alone. Both derivatives in
\cref{eq:force-torque} therefore vanish on the $j=i$ term.

The second term of \cref{eq:langevin-position} is self-propulsion at speed $v_0$ along
the rod axis. Here $k_i$ is a dichotomous noise: $k_i=\pm1$ with
equal probability, switching at rate $\frev$,
\begin{equation}
    \label{eq:k-dichotomous}
    \la k_i(t)k_j(t')\ra = \delta_{ij}\,\ee^{-2\frev|t-t'|}.
\end{equation}
This is the periodic reversal of \textit{Myxococcus xanthus} gliding motility \citeTODO{}. On timescales long compared with $\frev^{-1}$ we approximate $k_i$ by Gaussian white noise,
$\la k_i(t)k_j(t')\ra \approx \frev^{-1}\delta_{ij}\delta(t-t')$. The coarse-graining of
\cref{sec:model-coarse} uses It\^o calculus and therefore requires this approximation.
It is the only approximation made there.

\prompt{Maybe better to talk in terms of multiplicative noise}
Reversal does more than randomize the direction of self-propulsion. Because $k_i$ is
fast and multiplies a state-dependent vector, expanding \cref{eq:langevin-position} to
second order in $k_i\,\d t$, with $(k_i\,\d t)^2=\frev^{-1}\d t$, leaves a systematic
drift alongside the fluctuation (\cref{app:dean-microscopic}). After coarse-graining the
drift becomes the active force density in \cref{eq:model-rho}, and the fluctuation
becomes the density noise of the same equation.

\subsection{Coarse-graining that preserves fluctuations}
\label{sec:model-coarse}
% dean.tex 174-359 ("Dean's approach").
Dean's equation \cite{dean1996} is an \emph{exact} field-theoretic Langevin equation for
the density of interacting Langevin processes. Our model \eqref{eq:model} adds an
orientation, a self-propulsion and a reversal noise. The same derivation then gives an
exact equation for the joint density over the state space $(\r,\theta)$; the only
approximation is the Gaussian white-noise limit taken in \cref{sec:model-micro}.

Write the single-particle density $\varrho^{(i)}_t(\r,\theta) =
\delta(\r-\r_i(t))\,\delta(\theta-\theta_i(t))$, and let
$\varrho_t(\r,\theta)=\rho_0^{-1}\sum_i\varrho^{(i)}_t(\r,\theta)$ be the joint density
over $(\r,\theta)$, made dimensionless by $\rho_0$. The interaction is
pairwise, so the force and torque fields \cref{eq:force-torque} are gradients of a single
functional of $\varrho_t$,
\begin{equation}
    \label{eq:pair-free-energy}
    \calF[\varrho] = \frac{\rho_0^2}2\!\int\!\d\r\d\r'\,\d\theta\d\theta'\,
    V(\r\!-\!\r'\!,\theta,\theta')\,\varrho(\r,\theta)\,\varrho(\r'\!,\theta'),
\end{equation}
namely $\F=-\rho_0^{-1}\partial_\r\,\delta\calF/\delta\varrho$ and
$\Gamma=-\rho_0^{-1}\partial_\theta\,\delta\calF/\delta\varrho$. Applying It\^o's lemma
to $\varrho^{(i)}_t$ and summing over $i$ then gives a closed equation for $\varrho_t$
alone (\cref{app:dean-approach}),
\begin{equation}
    \label{eq:final-langevin}
    \begin{aligned}
    \partial_t\varrho_t(\r,\theta) = {}&
    \frac1{\xi_0}\pd{}{\r}\!\cdot\!\left(\varrho_t\pd{}{\r}\frac{\delta\calF}{\delta\varrho}\right)
    + \frac1{\xir}\pd{}{\theta}\!\left(\varrho_t\pd{}{\theta}\frac{\delta\calF}{\delta\varrho}\right)
    \\
    &+ \frac{2\zeta}{\xi_0}\,p_\alpha(\theta)p_\beta(\theta)\,\pd{^2\varrho_t}{r_\alpha\partial r_\beta}
    + \Lambda\,\pd{^2\varrho_t}{\theta^2}
    \\
    &- \pd{}{\r}\!\cdot\!\left(\sqrt{\varrho_t}\,\p(\theta)\,k(\r,\theta,t)\right)
    \\
    &- \pd{}{\theta}\!\left(\sqrt{\varrho_t}\,\lambda(\r,\theta,t)\right).
    \end{aligned}
\end{equation}
The last line of \cref{eq:final-langevin} carries two noise fields, $k(\r,\theta,t)$ and
$\lambda(\r,\theta,t)$. Both are Gaussian white noises in position, orientation and
time,
\begin{align}
    \label{eq:noise-correlators}
    \la k(\r,\theta,t)\,k(\r',\theta',t')\ra &= \frac{v_0^2}{\frev\rho_0}\;\Delta,
    \nonumber\\
    \la \lambda(\r,\theta,t)\,\lambda(\r',\theta',t')\ra &= \frac{2\Lambda}{\rho_0}\;\Delta,
\end{align}
with $\Delta\equiv\delta(t-t')\delta(\r-\r')\delta(\theta-\theta')$, and they are
uncorrelated with each other. Their amplitudes are inherited from the reversal noise
$k_i$ and the rotational noise $\lambda_i$ of \cref{eq:model}.
\Cref{app:dean-approach} derives that relation, together with the multiplicative factor
$\sqrt{\varrho_t}$ that accompanies both fields.

\placeholder{%
  Discussion of Dean's equation, to be written once the field equations are in place:
  \begin{itemize}
    \item \Cref{eq:final-langevin} is a Langevin equation for a stochastic field, not
      an evolution equation for an ensemble average with a fluctuation correction
      added afterwards. Say what that distinction buys.
    \item The noise is multiplicative, its amplitude set by $\varrho_t$ itself, and
      conservative in the translational channel.
    \item Nothing is added at the continuum level. The noise is what survives of the
      particle-level reversal noise of \cref{sec:model-micro}, and the only
      approximation between the two is the Gaussian white-noise limit.
    \item Contrast with the two routes that lose it: a mean-field continuum limit taken
      from the start, and positing continuum equations of motion directly, as is
      standard for active nematics \cite{marchetti2013hydrodynamics}. Both return a
      deterministic equation with no microscopic noise source left to carry through,
      and neither can be repaired afterwards by an equilibrium (FDT) argument,
      per \cref{sec:intro}.
    \item The amplitude in \cref{eq:noise-correlators} is $v_0^2/\frev$, built from the
      same two microscopic numbers, and in the same combination, as the systematic
      drift produced by the same expansion in \cref{sec:model-micro}. One measured
      biological parameter controls the mean active stress and the noise together.
  \end{itemize}}

\subsection{Free energy, closure, and the stochastic field equations}
\label{sec:model-closure}
The joint density $\varrho_t$ carries far more information than we need. The rest of the
paper works with its lowest angular moments, the fluctuating density and nematic order
parameter,
\begin{align}
  \label{eq:density-definition}
    \rho(\r,t) = \int\!\d\theta\,\varrho_t(\r,\theta),
    \quad
    \Q(\r,t) = \int\!\d\theta\,\varrho_t(\r,\theta)\,\hat\Q(\theta),
\end{align}
with $\hat\Q(\theta)=2\p(\theta)\p(\theta)-\I$ the single-rod nematic tensor.%
\footnote{%
The conventional nematic order parameter is normalized by the local density,
$\Q=\rho^{-1}\!\int\!\d\theta\,\varrho\hat\Q$, and is bounded by unity. We use the
unnormalized moment \cref{eq:density-definition} instead; the two coincide in the
incompressible limit $\rho\to1$ taken throughout (\cref{app:dean-normalisation}).}
In this subsection we reduce \cref{eq:final-langevin} to a closed pair of stochastic
equations for these two fields.

We do not know the interbacterial potential
${V(\r-\r',\theta,\theta')}$ in any detail, so we do not attempt to evaluate
\cref{eq:pair-free-energy} from it. Instead we assume that the free energy depends on
$\varrho$ only through the two moments \cref{eq:density-definition},
\begin{align}
    \label{eq:free-energy-parametrisation}
    \calF[\varrho] = F[\rho,\Q] = \int\!\d\r\;\mathfrak f\bigl(\rho(\r),\Q(\r)\bigr),
\end{align}
which is both a physical assumption and a choice of parametrization. It gives a chain
rule,
\begin{align}
    \label{eq:chain-rule}
    \frac{\delta\calF}{\delta\varrho(\r,\theta)}
    = \frac{\delta F}{\delta\rho(\r)}
    + \frac{\delta F}{\delta\Q(\r)}\!:\!\hat\Q(\theta),
\end{align}
This truncation of the $\theta$ dependence to a single harmonic is what makes
\cref{eq:final-langevin} integrable over $\theta$ (\cref{app:dean-free-energy}).

For free energy density$\mathfrak f$ we take the sum of an incompressibility penalty and a Landau-de Gennes
nematic term,
\begin{align}
    \label{eq:free-energy-choice}
    \mathfrak f(\rho,\Q) = {}& \frac B2(\rho-1)^2 + \frac a2(\Q\!:\!\Q)
    + \frac b4(\Q\!:\!\Q)^2
    \nonumber\\
    &+ \frac K2(\partial_\r\Q)\!:\!(\partial_\r\Q).
\end{align}
The first term enforces near-incompressibility of the monolayer with a large bulk
modulus $B$; the rest is the standard nematic free energy, with $a<0$ favoring order and $K$ penalizing distortion. The single $K$ is the one-constant approximation of the Frank free energy, in which the splay and bend moduli of the two-dimensional nematic are taken equal \citeTODO{}.

Integrating \cref{eq:final-langevin} over $\theta$ and keeping the leading terms in
$1/B$ (\cref{app:dean-angle}) gives a continuity equation for the density,
\begin{align}
  \label{eq:model-rho}
  \partial_t\rho = -\nabla\!\cdot\!\Bigl[
    \frac1{\xi_0}\bigl(-B\nabla\rho + \Q\!:\!\nabla\H + \fa\bigr)
    - \bEta_{\rho}(\r,t)\Bigr].
\end{align}
The first two terms in the bracket are passive, inherited from the pair potential through
the free energy. The term $-B\nabla\rho$ resists compression, and $\Q\!:\!\nabla\H$
carries cells along gradients of Nematic force field 
$\H(\r)=-\delta F/\delta\Q= -a\Q - b(\Q\!:\!\Q)\Q + K\nabla^2\Q$ \citeTODO{}. The third term is the active force density
\begin{align}
  \label{eq:active-force}
  f^a_\alpha(\r) = -\zeta\,\partial_\beta Q_\ab(\r)
  \quad\text{with}\quad
  \zeta = \frac{\xi_0v_0^2}{4\frev},
\end{align}
the divergence of the active nematic stress $-\zeta\Q$. Positive $\zeta$ denotes extensile active stress.
The coefficient comes entirely from self-propulsion and reversal, so fast reversin colonies are less active.

The last term is the noise of the density flux. Its autocorrelation follows from
\cref{eq:noise-correlators},
\begin{multline}
  \label{eq:eta-rho-autocorrelation}
  \la \eta^{\rho}_\alpha(\r,t) \,\eta^{\rho}_\beta(\r',t') \ra =
  \\
  =
  \frac{2\zeta}{\xi_0\rho_0}
  \left[Q_\ab(\r)+\rho(\r) \delta_\ab\right]
 \delta(\r-\r')\delta(t-t') ,
\end{multline}
an isotropic piece proportional to $\rho$ and an anisotropic piece proportional to $\Q$.
Both scale with $\zeta$ of \cref{eq:active-force}, so the noise strength is set by the
same reversal rate $\frev$ and self-propulsion speed $v_0$ as the active force. The
rotational noise does not appear: every term it produces is a $\theta$ derivative, and
$\int\!\d\theta\,\partial_\theta(\cdot)=0$. Noise $\bEta_{\rho}$ sits inside the divergence, so it
conserves the number of cells \prompt{the word"exactly" is confusing in this context. what you mean by that?} exactly and not only on average.

\Cref{eq:model-rho} is a continuity equation,
$\partial_t\rho=-\nabla\!\cdot\!\bs J$. If we write $\bs J = \rho \v$, where $\v(\r,t)$
is the velocity field, and set $\rho\approx1$ in the incompressible limit, we get a force balance
\begin{align}
  \label{eq:velocity}
  \xi_0 \v = -B\nabla\rho + \Q\!:\!\nabla\H + \fa - \xi_0 \bEta_{\rho}.
\end{align}
where the friction force density $\xi_0 \v$ on the cell is balanced by the passive interactions, active nematic force and the noise. $B\nabla\rho$ acts like a pressure gradient.
The elastic force $\Q\!:\!\nabla\H$ redistributes
cells to relieve orientational stress. Only $\fa$ does not derive from Free energy $F$.Since $\fa\propto\nabla\!\cdot\!\Q$, it vanishes in a
uniform texture and is largest where the texture bends fastest, at a defect core
(\cref{sec:defect-steady}).

We also need a dynamical equation for the nematic order parameter $\Q$. Multiplying
\cref{eq:final-langevin} by $\hat\Q(\theta)$ and integrating over $\theta$ does not close
the system: the angular averages like $\la\hat\Q\hat\Q\ra_\varrho$ appear, and $\rho$ and $\Q$ do not determine them. We use
the lowest-order maximum-entropy closure,
$\langle\hat\Q\hat\Q\rangle_\varrho \simeq \rho\,\J$
with $\J$ the rank-four identity on symmetric traceless tensors. It is accurate to
$O(S^2)$ in the strength of the nematic order $S=|\Q|$. \Cref{app:closure} derives the
higher-order terms, and the limitations of this closure, which affect only the
numerical solutions.

% \prompt{I desiced to split noise term into two parts: one conserved noise, the other non-conserved noise. eta^Q is correlated with eta^rho, while f^Q is uncorrelated with either. Worth mentioning that Eta^Q is rank-3 tensor. This notation should also be used in the appendix, which is different then the dean.tex. Update the rest of the section accordingly.}
With the closure in place we obtain
\begin{align}
  \label{eq:model-Q}
  \partial_t \Q \!=\! -\frac{1}{\xi_0} \nabla\!\cdot\!\Bigl(\!-\Q B \nabla\rho
  +\! \nabla \H-\!\bEta^\Q\Bigr)\!+\! \frac4{\xir} \Ha\! + \f^Q.
\end{align}
The first two terms inside the divergence are the Onsager conjugates of the passive
terms of \cref{eq:model-rho}. The third is the
nematic force field with coefficients renormalized by the activity,
\begin{align}
  \label{eq:H-active}
  \Ha &= -a'\Q - b(\Q\!:\!\Q)\Q + K'\nabla^2\Q,
  \nonumber\\
  a' &= a+\Lambda\xir, \qquad K' = K+{\zeta\xir}/({4\xi_0}).
\end{align}
It is $a'$ and $K'$, not the bare $a$ and $K$, that set the defect profile solved for in
\cref{sec:defect-steady}.

The local and the flux terms of \cref{eq:model-Q} do different jobs. The local term
$\tfrac4{\xir}\Ha$ relaxes $\Q$ toward a minimiser of the effective free energy\footnote{\prompt{H is not fully local, it depends on the laplacian of Q. what we mean is the non-trasport term. Maybe you can use better naming to avoid this footnote}}, and
fixes the shape of the defect. The flux terms leave that shape almost unchanged, but
through $\Q B\nabla\rho$ they couple the texture to the density field and translate the
defect. \Cref{sec:defect-steady} solves for that translation.

The noise $\f^Q$ is given in \cref{app:dean-angle}. Unlike $\bEta_\rho$ it carries both
the reversal noise $k(\r,\theta,t)$ and the rotational noise $\lambda(\r,\theta,t)$, and
it is correlated with $\bEta_\rho$, since the two share the same $k$.

In the rest of the paper we drop the terms proportional to $\nabla\H$ from
\cref{eq:model-rho,eq:model-Q}. They introduce fourth-order spatial derivatives of $\Q$,
which we cannot treat analytically and numerically. The two terms must
go together, because they are Onsager conjugates: dropping one alone would break
reciprocity. They are also the smaller of the two channels by which orientational order
relaxes. Rods can turn in place, at a rate set by the rotational friction $\xir$, or rods
of different orientation can be exchanged between neighbouring regions, at a rate set by
the translational friction $\xi_0$. The second channel is a flux divergence and costs two
extra powers of wavevector, so the ratio of the two rates is $(k\ell_\star)^2$ with
$\ell_\star=\sqrt{\xir/\xi_0}=0.14~\micron$ (\cref{app:parameters-choice}). At the defect
core this suppresses the discarded channel by
$(\ell_\star/\elld)^2\simeq2\times10^{-2}$.

\todo{We need to understand what physics we are missing by discarding these terms.}

\Cref{eq:model-rho,eq:model-Q} are the fluctuating active nematic equations this paper
solves. \prompt{activity enters also in $\Ha$} The activity enters them twice under one coefficient: as the force $\fa$ and as
the amplitude of $\bEta_\rho$, both fixed by $\zeta$ and hence by the reversal rate. The
density and the nematic field are coupled in both directions, through
$\fa\propto\nabla\!\cdot\!\Q$ in one equation and through $\Q B\nabla\rho$ in the other.
That coupling sets a $+1/2$ defect in motion in \cref{sec:defect-steady}.


\prompt{Equation for first moment $\int\!\d\theta\,\varrho_t\,\p(\theta)$ is meaningless. Note that this is not observable, The polarity of the cell is the $p(\theta_i) k_i$, i.e. multiplied by the noise. Additionally, the first moment as defined does not appear in the closure procedure. So unnecesary and meaningless.}
There is no third equation, for the polar order $\int\!\d\theta\,\varrho_t\,\p(\theta)$.
Reversal removes it: $k_i$ has zero mean, so self-propulsion survives the coarse-graining
only at second order, and the moment hierarchy couples $\rho$ directly to $\Q$. The
monolayer is a nematic, not a polar active fluid.

Everything that follows comes from solving \cref{eq:model-rho,eq:model-Q} at a $+1/2$
defect, with the parameters collected in \cref{app:parameters}.

% ============================================================
\section{Fluctuations around a $+1/2$ defect}
\label{sec:defect}
% structure.md restructure, 2026-09-21. dean.tex 681-1266.
% ============================================================
Solving \cref{eq:model-rho,eq:model-Q} at a $+1/2$ defect takes three steps: find the
 steady state of the noiseless equations
(\cref{sec:defect-steady}); linearize the stochastic equations around it; and solve the
resulting linear field theory for the equal-time covariance of the fluctuations
(\cref{sec:defect-fluct}).

We solve \cref{eq:model-rho,eq:model-Q} numerically in
\cref{app:comoving,app:lyapunov}. The main text gives analytical solutions of simplified
equations, which agree well with the numerics. The analytical solutions also show how the
result depends on the parameters of the model, which the numerical solutions do not.

\subsection{The steady state}
\label{sec:defect-steady}
% Pages ~1.5. dean.tex 705-747, 819-894, 2460-2597 (solver moved to appendix, results summarised here).

We first solve a simplified system. The $+1/2$ defect moves because the density and the
nematic order are cross-coupled, so we drop the cross-coupling term in the $\Q$ equation:
\begin{align}
  \label{simplified-noiseless}
  0 = \nabla\cdot\left[B\nabla\rho - \fa\right] \qquad \Ha = 0
\end{align}
The solution of the second equation is given by Pad\'{e} approximation (see Appendix ??) \citeTODO{}:
\begin{align}
  \label{eq:pade-approximation}
  \Q(r,\phi) = S(r)\hat\Q(\phi/2)
  \qquad
  S(r) = S_0 f_\pade(r/\elld)
  \\
  f_\pade(x) = x \sqrt{\frac{0.34 + 0.07 x^2}{1 + 0.41 x^2+0.07 x^4}}
\end{align}
Two scales appear. The far-field nematic order $S_0=\sqrt{-a'/(2b)}$ is the value $S(r)$
saturates to, and the core size $\elld=\sqrt{K'/(-a')}$ is the distance over which $S$
recovers from $S(0)=0$ at the defect. Both are built from the activity-renormalized $a'$
and $K'$ of \cref{eq:H-active}, not from the bare coefficients of the free energy
density \eqref{eq:free-energy-choice}. \Cref{app:parameters-choice} gives their values and \cref{app:comoving} the
accuracy of the approximant \cite{copenhagen2021,han2025local}.

Now, the first equation in \cref{simplified-noiseless} is a Poisson equation for the density
with an exact solution:
\begin{align}
  \label{eq:rho-analytic}
  \rho(r,\phi) = 1 - \frac{\zeta}{B}\left[ S(r) - \frac{r}{R}S(R)\right]\cos \phi
\end{align}
where ($r$, $\phi$) are the polar coordinates and we impose $\rho(R,\phi) = 1$ on the boundary $r=R$.
The density therefore departs from uniformity only at order $\zeta/B$.

\prompt{below in this section, it is mentioned 3 times at least that solution is in the comoving frame. Clearly, the section is not well written. I feel that same thing is told multiple times in different way, every time adding small new detail. This is bad, rewrite it}
\Cref{app:comoving} obtains the numerical solution of the noiseless
\cref{eq:model-rho,eq:model-Q} without either simplification. Because of the
cross-coupling between density and nematic order the defect moves, so that solution is
obtained in a frame comoving with the core. \Cref{fig:rho1d-dirichlet} shows the section
of the density field at $y=0$ against \cref{eq:rho-analytic}, and
\cref{fig:rhodensity,fig:rho3d} the full two-dimensional field.

The defect moves \emph{opposite} to the active force. The active force does not push the
defect; it pushes the cells. The cells build a density gradient until the diffusive flux
cancels the active one, $B\nabla\rho\simeq\fa$. The nematic order parameter, however,
couples to the density gradient and not to the active force. The nematic flux in
\cref{eq:model-Q} is $\bs J_\Q = -\Q B\nabla\rho/\xi_0 \simeq -\Q\fa/\xi_0$, so $\Q$ is
transported with velocity $-B\nabla\rho/\xi_0$, down the density gradient and therefore
against $\fa$.

For the $+1/2$ defect \cref{eq:pade-approximation} gives
\begin{align}
  \label{eq:active-force-defect}
  f^\text{a}_x(r, \phi) = -\zeta \left(S'(r) +S(r)/r\right),
  \qquad f^\text{a}_y = 0,
\end{align}
a vector field along the defect axis with no transverse component. It points along $-x$,
the side on which the director bends around the core (\cref{fig:rhodensity}), so the
defect drifts along $+x$ at a speed $u_x = 0.20~\micron/\minute$
(\cref{app:comoving-results}).

The numerical solution is shown in the comoving frame, with the defect core at the
origin. \Cref{eq:rho-analytic} is exactly antisymmetric under $x\to-x$; the
numerical solution is not, and that asymmetry, visible in \cref{fig:rho1d-dirichlet}, is
the signature of the drift. The density coupling terms in the $\Q$ equation also deform
the defect slightly. \Cref{fig:strength,fig:defectangle} show the deviation from the
Pad\'e approximation \eqref{eq:pade-approximation}: \prompt{In the main text I just want to say that deviation is small, about 1 degree. More details in the figure caption}the director rotates by at most one
degree out to $r=5~\micron$, and the deviation grows with radius.
\Cref{fig:activeforce} shows the same comparison for the active force
\cref{eq:active-force-defect}.


\begin{figure}[tbp]
  \centering
  \includegraphics[width=0.85\linewidth]{rho1d_Dirichlet.png}
  \caption{Steady-state density along the defect axis, $(\rho-1)/(\zeta/B)$ at $y=0$.
    Solid: the comoving-frame numerical solution of \cref{eq:model-rho,eq:model-Q}
    (\cref{app:comoving}), reservoir (Dirichlet) boundary condition. Dashed:
    \cref{eq:rho-analytic}, the closed-form solution of the simplified system, which is
    antisymmetric under $x\to-x$. The density is enhanced behind the defect and depleted
    ahead of it, and the departure from antisymmetry is the effect of the drift.}
  \label{fig:rho1d-dirichlet}
\end{figure}

\begin{figure}[tbp]
  \centering
  \includegraphics[width=0.85\linewidth]{rhodensity.png}
  \caption{Steady-state density field $(\rho-1)/(\zeta/B)$ around the $+1/2$ defect,
    comoving frame, with the core at the origin. Dark curves are the nematic director.
    The director bends around the core on the $-x$ side and splays along the axis on the
    $+x$ side; the density is enhanced on the former and depleted on the latter.
    \Cref{fig:rho3d} shows the same field as a surface.}
  \label{fig:rhodensity}
\end{figure}

\begin{figure}[tbp]
  \centering
  \includegraphics[width=0.85\linewidth]{rho3d.png}
  \caption{Converged $+1/2$ defect steady-state density field $\rho_\ss(\r)$,
    comoving-frame numerical solution (\cref{sec:defect-steady}), Dirichlet
    (reservoir) boundary condition. The core is pinned at the origin by construction
    (\cref{app:comoving}); the sealed-wall solution is visually indistinguishable at
    this scale, since the two boundary conditions differ in the far field, and hence in
    $\u$, not in the core structure shown here. \Cref{fig:rhodensity} shows the same
    field as a plan view.}
  \label{fig:rho3d}
\end{figure}


Because the defect drifts, there is no stationary solution in the laboratory frame, and
we solve numerically in a frame moving with the defect. \Cref{app:comoving} gives the
comoving-frame equations, the constraint that fixes the frame velocity $\u$, and the
boundary conditions.

The closed-form density solution is a check on the numerics. Fixing $\Q$ to the Pad\'e
profile \cref{eq:pade-approximation} is the same simplification made in
\cref{simplified-noiseless}, so the analytic $\rho$ and the converged numerical
$\rho_\ss$ are independent routes to the same field. They agree: both are depleted ahead
of the defect and enhanced behind it, with amplitude $O(\zeta/B)$. The density profile
enters only the \emph{mean} pressure, which \cref{sec:nucleation-theory} shows to be far
below the fluctuations, so this agreement matters for consistency rather than for the
final result.


\subsection{Linearization and the fluctuation spectrum}
\label{sec:defect-fluct}
% Pages ~2.5. dean.tex 748-818, 895-1266.
Let $\delta\rho(\r,t)=\rho(\r,t)-\rho_\ss(\r)$ and $\delta\Q(\r,t)=\Q(\r,t)-\Q_\ss(\r)$.
Linearizing the comoving-frame form of \cref{eq:model-rho,eq:model-Q} around the
steady state of \cref{sec:defect-steady} gives
(\cref{app:dean-angle})
\prompt{check the newly added noise term, update the text accordingly}
\begin{subequations}
    \label{eq:linearized}
\begin{align}
    &\partial_t\delta\rho =-\nabla\!\cdot\!\Bigl[
    \frac1{\xi_0}\bigl(-B\nabla\delta\rho -\zeta\,\nabla\!\cdot\!\delta\Q \bigr) + \bEta_\rho(\r,t)\Bigr],
    \label{eq:linearized-rho}
    \\
    &\partial_t\delta\Q = \frac B{\xi_0}\nabla\!\cdot\!(\Q_\ss\nabla\delta\rho)
    + \nabla\!\cdot\!(\delta\Q\,\bs w_\ss) + \frac4{\xir}\delta\Ha
    \nonumber\\
    &\qquad\qquad
    +\frac{1}{\xi_0}\nabla\cdot\bEta^Q(\r,t)+ \f^Q(\r,t),
    \label{eq:linearized-Q}
\end{align}
\end{subequations}
with $\bs w_\ss\equiv\u+(B/\xi_0)\nabla\rho_\ss$ the steady advection velocity and
$\delta\Ha$ the linearization of \cref{eq:H-active} at $\Q_\ss$. This is a linear
stochastic field theory on the background $(\rho_\ss,\Q_\ss)$. \Cref{app:lyapunov}
obtains the equal-time covariance of $\delta\rho$ and $\delta\Q$ numerically, by solving
a Lyapunov equation.

The main text presents the analytically tractable case, in which only the density
fluctuates and $\delta\Q=0$. This is the appropriate limit when the fluctuations of the
nematic order parameter are small compared with the noise in \cref{eq:linearized-rho}.
The dominant source of fluctuations is indeed the reversals of the bacteria
\cite{han2025local}.

The defect then enters \cref{eq:linearized-rho} in exactly one place, through the
$\Q_\ss$ term in the noise autocorrelation \cref{eq:eta-rho-autocorrelation}. The defect
shape is felt by the fluctuations and not by the steady state.

\prompt{This is still bad, we introduce the cutoff before talking about it. It is wrong to call eq:rho-fourier as a fourier transform, as it holds the cutoff. We talk about it too far away from the introduction of the cutoff, which is confusing.}
Taking the Fourier transform of \cref{eq:linearized-rho} over space and time gives
\begin{align}
  \label{eq:rho-fourier}
  \delta\tilde\rho(\k,\omega) =
  \frac{\ii\,\k\cdot\tilde \bEta_\rho(\k,\omega)}{\ii\omega+\gamma_0 k^2}\,\ee^{-k^2\elln^2/4},
  \qquad \gamma_0 \equiv B/\xi_0,
\end{align}
with $\k$ the wavevector and $k=|\k|$. The noise is conservative, so its power carries a
factor $k^2$ that cancels the $k^2$ of the relaxation rate $\gamma_0k^2$. The spectrum is
then flat in $\k$, and the equal-time variance
$\la\delta\rho(\r,t)\delta\rho(\r',t)\ra_\ss$ diverges as $\r'\to\r$. The model has no
short-length cutoff of its own, so we impose the physical one, the width of a bacterium, which is the shortest length scale of our system.
\begin{equation}
    \label{eq:cutoff}
    \elln = 0.5~\micron.
\end{equation}
It enters \cref{eq:rho-fourier} as the exponential suppression of wavevectors above
$1/\elln$.

Consider first the isotropic case $\Q_\ss=0$. Inverting \cref{eq:rho-fourier} over space and time gives
\begin{align}
  \label{eq:rho-autocorrelation-isotropic}
  \la\delta\rho(\r,t)\delta\rho(\r',t')\ra =
  \frac{\zeta}{2\pi B\rho_0 \elln^2}
  \frac{
      \exp\!\!\left(\!-\frac{|\r-\r'|^2}{2 \elln^2+4 \gamma_0 |t-t'|}\right)
      }{1 + 2 \gamma_0 |t-t'| / \elln^2}
\end{align}
so that the equal-point, equal-time variance is the constant pedestal
\begin{align}
  \label{eq:pedestal}
  \sigmarho(\r) = \sigmarho\big|_\text{pedestal}
  = \frac{\zeta}{2\pi B\rho_0\elln^2},
\end{align}
\Cref{eq:rho-autocorrelation-isotropic} also gives the two scales over which a density
fluctuation stays coherent: a correlation length equal to the cutoff $\elln$, as it must
be once the spectrum is flat below it, and a correlation time
\begin{align}
 \label{eq:tau-c}
 \tau_c = \frac{\elln^2}{2 \gamma_0},
\end{align}
the time a fluctuation of size $\elln$ takes to diffuse away. Neither scale involves the
defect, and we assume the anisotropy does not change them much. \Cref{sec:nucleation-results}
uses $\tau_c$ to compute the layering probability.

Restoring the anisotropic piece of \cref{eq:eta-rho-autocorrelation}, the one that
carries $\Q_\ss$, gives the total variance (\cref{app:lyapunov-finite-corr})
\begin{align}
  \label{eq:total-gauss}
  \sigmarho(\r) = \sigmarho\big|_\text{pedestal}
  \left(1 + \frac{\elln^2}{4}\bigl(\kappa\star\partial_i\partial_jQ^\ss_{ij}\bigr)(\r)\right),
\end{align}
a convolution of the double divergence of the steady nematic texture against the
isotropic real-space kernel
\begin{equation}
  \label{eq:kappa-real}
  \kappa(w) = \frac{1-\ee^{-2w^2/\elln^2}}{2\pi\,w^2}.
\end{equation}
\Cref{eq:total-gauss} is the central result of \cref{sec:defect}. 
\prompt{The following sentence is not good. the shot noise is wrong terminology. Ideally, I would like to see something like: Defect modulates around the uniform pedestal, also mention filtering etc. Think about it}
The pedestal is uniform: it is the price of resolving conserved shot noise at scale $\elln$, and it carries no information about the defect.
The defect enters only through the bracket, as
the double divergence of $\Q_\ss$ blurred by $\kappa$. Beyond the cutoff $\kappa$ decays
as $1/2\pi w^2$, so nematic structure on every scale between $\elln$ and the core
contributes to the variance at a given point. Inside the cutoff the exponential flattens
$\kappa$ to $1/\pi\elln^2$, so structure finer than a cell contributes nothing.

Note what does \emph{not} appear in \cref{eq:total-gauss}. The steady density $\rho_\ss$
drops out beyond its far-field value, so neither the density profile of
\cref{sec:defect-steady} \st{nor the drift speed $u$ plays any part} \prompt{do not mention the drift speed $u$ in the main text}. The defect is visible in
the variance only through $\partial_i\partial_jQ^\ss_{ij}$, the same combination that
carries the active force in \cref{eq:active-force}.

\Cref{eq:total-gauss} is exact at leading order in $1/B$ and rests on the $\delta\Q=0$
truncation made above. \Cref{fig:variance} compares it against the full numerical
solution of \cref{eq:linearized}, retaining $\delta\Q\ne0$. Along the defect axis the
modulation is odd in $x$ and vanishes at the core, because
$\partial_i\partial_jQ^\ss_{ij}=\partial_x[S'(r)+S(r)/r]$ is odd. The variance is raised
on the $-x$ side, where it peaks about $1.5~\micron$ from the core at
$\nu_0\equiv[\sigmarho(\r_\mathrm{max})/\sigmarho(\infty)]-1=3.1\%$ for
$\elln=0.5~\micron$, and lowered by about as much on the $+x$ side. Across the axis the
modulation is an order of magnitude smaller.
$\nu_0$ decreases with decreasing $\elln$ as shown in \cref{fig:ln_curve} and eventually reaches 0 at $\elln=0$, meaning that without the cutoff, there would be no modulation at all.
Three percent is the entire signal, and in \cref{sec:nucleation} we show that 
this small modulation can still lead to a larger modulation of the layering probability around the defect.

\begin{figure}[tbp]
  \centering
  \includegraphics[width=\linewidth]{varrho.png}
  \caption{Equal-point density variance around a $+1/2$ defect, normalized by the
    pedestal of \cref{eq:pedestal}, along the defect axis ($\xi=x$, $y=0$) and across it
    ($x=0$, $\xi=y$). Dashed: the analytic result at $\delta\Q=0$,
    \cref{eq:total-gauss}. Solid: the numerical Lyapunov solution retaining
    $\delta\Q\ne0$ (\cref{app:lyapunov}). Along the axis the modulation is odd and
    vanishes at the core, reaching $+\nu_0=3.1\%$ on the $-x$ side and about $-3\%$ on
    the $+x$ side, both roughly $1.5~\micron$ out, for $\elln=0.5~\micron$. Across the
    axis it is much weaker.}
  \label{fig:variance}
\end{figure}

% ============================================================
\section{Layer formation as a nucleation event}
\label{sec:nucleation}
% Was "Extrusion and layer formation" (Results-2D/2E in v2). Renamed and
% resectioned per structure.md restructure, 2026-09-21.
% Pages ~2.5 (compact by design; full ranges/sensitivity -> Appendix F.3).
% dean.tex 1281-1487 ("Extrusion rate": Virtual work, Parameter estimates,
% The nucleation threshold, Layering probability).
% ============================================================

\subsection{Energetics of layer formation}
\label{sec:nucleation-energetics}
% dean.tex 1284-1360 (main-text portion only; full ranges to Appendix F.3).

We treat the monolayer as infinite, adhering to the substrate and to itself, and wetted
by a thin film of water. A patch of cells of area $A$ and perimeter $\ell$ lifts out of
it against three costs and with two gains. The costs: the footprint of area $A$ is peeled
off the substrate at $\Ws$ per unit area; the side wall of area $\ell h$, for a layer of
thickness $h$, was in-plane cell-cell contact and is now free, at $\Wc$ per unit area;
and the wetting film must re-coat that side wall, growing the water-air interface by
$\ell h$ at the surface tension $\gammaw$ of water. The gains: the underside of area $A$
lands on the cells beneath and makes fresh cell-cell contact, and the compressed parent
monolayer relaxes. Collecting the five terms,
\begin{equation}
    \label{eq:virtual-work}
    \Delta E = \Ws A + \Wc\ell h + \gammaw\ell h - \Wc A + \Delta E_e,
\end{equation}
with $\Delta E_e$ the change in elastic strain energy of the monolayer.

The cell-cell adhesion $\Wc$ enters \emph{twice} with opposite signs: as a perimeter cost
on the side faces and as an area gain on the bottom face. The substrate adhesion $\Ws$ is
a pure area cost and the surface tension $\gammaw$ a pure perimeter cost. Every perimeter
term is therefore positive while the area terms are mixed in sign, and this asymmetry is
the origin of the nucleation barrier in \cref{sec:nucleation-theory}. None of the three
energies has been measured for \textit{M.\ xanthus}. $\gammaw$ exceeds any plausible
adhesion energy by an order of magnitude, so it is the wetting film and not the adhesion
that holds the monolayer flat.

The elastic term is the one that drives the transition. A region compressed to density
$\rho$ above the relaxed value carries a bulk strain $\epsilon=\rho-1$ in the units of
\cref{eq:free-energy-choice} and stores energy $\tfrac B2\epsilon^2A_d$ over the
deformed area $A_d$, with $B$ the same two-dimensional bulk modulus that sets the
incompressibility penalty in the field theory. Extruding $A$ worth of cells dilutes that
region, which holds its area while $\rho_0(1+\epsilon)A$ cells leave it, so its strain
falls to $\epsilon_f=\epsilon-(1+\epsilon)A/A_d$; the extruded cells themselves are
unconfined in the new layer and relax to zero strain. Hence
$\Delta E_e=\tfrac B2(\epsilon_f^2-\epsilon^2)A_d$, or
\begin{equation}
    \label{eq:dEe-full}
    \Delta E_e = \frac B2\left[-2\epsilon - 2\epsilon^2
    + (1+\epsilon)^2\frac A{A_d}\right]A,
\end{equation}
which to leading order in the small compression $\epsilon$ and the small area ratio
$A/A_d$ collapses to
\begin{equation}
    \label{eq:dEe}
    \Delta E_e \simeq -B\epsilon A = -pA,
\end{equation}
the work done by the local monolayer pressure $p=B\epsilon$ against the area it opens up.
The identification $p=B\,\delta\rho$ connects this section to \cref{sec:defect-fluct}:
the pressure in \cref{eq:dEe} is the field whose variance was computed there.

Two idealizations bound the result rather than change its structure
(\cref{app:extrusion-idealizations}). First, we assumed that the parent layer holds its
area fixed and dilutes. If instead it contracts to hold its density fixed, it loses
substrate contact both under the nucleus and at its own periphery, so the area peeled off
is $2A$ rather than $A$, while the work done against the confining pressure is unchanged.
The two limits differ by a rescaling of $\Ws$, not by a new term. Second, complete
wetting is the \emph{maximum} interfacial cost. If the water beads into a droplet instead
of coating the colony, less than $\ell h$ of new water-air interface is created and the
dominant term is reduced. This lowers the effective $\gammaw$ below the clean-water
value, independently of the surfactant argument below.

\Cref{tab:used-values} lists the values used here;
\cref{app:extrusion} gives the full literature ranges and a one-at-a-time
sensitivity analysis.

\begin{table}[tbp]
\centering
\begin{tabular}{c c l}
\hline
Parameter & Used value & Basis \\
\hline
$\gammaw$ & $6.5\times10^4~\Pa\cdot\micron$ & near clean air-water value \\
$\Wc$ & $0$ & no measured cell-cell cohesion \\
$\Ws$ & $2\times10^3~\Pa\cdot\micron$ & bacterial adhesion \\
$h$ & $0.5~\micron$ & cell diameter \\
$\Acell$ & $3~\micron^2$ & $\simeq1/\rho_0$ \\
$\elln$ & $0.5~\micron$ & cell width, \cref{eq:cutoff} \\
$B$ & $2\times10^5~\Pa\cdot\micron$ & fitted, \cref{sec:nucleation-results} \\
\hline
\end{tabular}
\caption{Parameters entering \cref{eq:virtual-work,eq:pcrit}, in units of
  $\Pa\cdot\micron=10^{-6}~\mathrm{J/m^2}$. Only the combinations $(\gammaw+\Wc)$ and
  $(\Ws-\Wc)$ appear in the results; full ranges and origins in
  \cref{app:parameters-extrusion}.}
\label{tab:used-values}
\end{table}

$\gammaw$ is set within $10\%$ of the clean air-water value. \textit{M.\ xanthus} excretes
surfactant, which would normally lower an interfacial tension substantially. The
criterion of \cref{sec:nucleation-theory}, however, requires only that the barrier
\emph{maximum} fall to one cell, not that the barrier vanish. That is a much weaker
requirement, and it leaves a near-water tension plausible. $\Wc=0$ is not a fit
but an experimental statement: \textit{M.\ xanthus} cells show essentially no cohesive
interaction, and the monolayer is held together by packing against the substrate. $B$ is
a model parameter, the stiffness of the incompressibility penalty in
\cref{eq:free-energy-choice}, not a measured material constant;
\cref{sec:nucleation-results} fits it to the comparison with experiment. Every other
entry, $\elln$ included, is fixed beforehand.

\subsection{Nucleation theory with a fluctuating barrier}
\label{sec:nucleation-theory}
% Merges dean.tex 1361-1390 ("The nucleation threshold") with the pressure-statistics
% opening of dean.tex 1392-1487 ("Layering probability") -- the deterministic threshold
% and the fluctuation statistics it is measured against belong in one subsection, since
% the section title now names the fluctuating barrier explicitly.
\Cref{eq:virtual-work} has the structure of classical nucleation theory once written as
a function of the nucleus area: a line cost growing as the perimeter and a bulk gain
growing as the area, with $\Delta E_e=-pA$ from \cref{eq:dEe} and $\ell=\sqrt{4\pi A}$
for a circular patch,
\begin{equation}
    \label{eq:dE-of-A}
    \Delta E(A) = (\Ws-\Wc-p)\,A + (\Wc+\gammaw)\,h\sqrt{4\pi A}.
\end{equation}
This rises as $\sqrt A$ at small $A$, turns over, and falls once the area term
dominates; the turning point is the classical critical nucleus, at radius
$r^*=(\Wc+\gammaw)h/(p+\Wc-\Ws)$.

The nucleation condition comes from the discreteness of the system. No nucleus can be
smaller than one cell, so \cref{eq:dE-of-A} is meaningful only down to $A\simeq\Acell$.
Once the \emph{maximum} of $\Delta E(A)$ has descended to that size, growth from the
smallest object that can exist is already downhill, and single-cell perturbations such as
thermal motion, substrate roughness or a division event complete the extrusion.
Equivalently, $r^*$ has shrunk to the size of one cell. The fluctuating quantity is then
not the size of the nucleus but the height and position of the barrier. The nucleation
condition is therefore
\begin{equation}
    \label{eq:nucleation-condition}
    \Delta E'(A)\big|_{A=\Acell} = 0,
\end{equation}
This condition is what makes a near-water surface tension compatible with observable
extrusion. Differentiating \cref{eq:dE-of-A} and solving gives the critical pressure
\begin{equation}
    \label{eq:pcrit}
    \pcrit = (\Ws-\Wc) + \sqrt{\frac{\pi h^2}{\Acell}}\,(\Wc+\gammaw),
\end{equation}
which is $r^*=\sqrt{\Acell/\pi}$ read backwards, and with \cref{tab:used-values} gives
$\pcrit=3.5\times10^4~\Pa\cdot\micron$. Note what $\pcrit$ does \emph{not} contain: no
$B$, no $\elln$, no property of the fluctuations. Adhesion, surface tension and cell
geometry fix it entirely. Everything from \cref{sec:defect-fluct} enters through the
statistics of $p$, which we turn to now.

The local pressure follows from the free energy density as $p=\rho\mathfrak f'-
\mathfrak f$, which for \cref{eq:free-energy-choice} and $\rho$ near unity is
$p\simeq B(\rho-1) = B\,\delta\rho$. Since $\delta\rho$ is Gaussian with the variance of
\cref{eq:total-gauss}, $p$ is Gaussian too, with
\begin{equation}
    \label{eq:sigma-p}
    \sigma_p^2(\r) = B^2\sigmarho(\r),
    \qquad
    \sigma_{p,\infty}^2 = \frac{B\zeta}{2\pi\rho_0\elln^2},
\end{equation}
the second expression its far-field value, from the isotropic pedestal of
\cref{eq:total-gauss}.

Only the width of that Gaussian matters, not its mean. The mean is the steady-state
pressure $B(\rho_\ss-1)$, the density profile of \cref{sec:defect-steady} read through
the same equation of state. That profile is of relative size $O(\zeta/B)$, so the mean
pressure is of order $\zeta$, a few $\Pa\cdot\micron$, against
$\sigma_{p,\infty}=1.8\times10^3~\Pa\cdot\micron$. The average compression at the defect
is smaller than the threshold by more than two orders of magnitude, and from here on we
take the pressure to be a zero-mean Gaussian of width $\sigma_p(\r)$.

$\pcrit$ of \cref{eq:pcrit} is a fixed number, set by adhesion and geometry, but the
pressure that must exceed it is a stochastic field whose variance the defect modulates
through \cref{eq:total-gauss}. This is the sense in which the barrier fluctuates.
Extrusion happens where and when a fluctuation crosses the fixed threshold, and the
probability of that at a point $\r$ is the Gaussian tail
\begin{align}
    \label{eq:p-layer}
    p_\text{layer}(\r) &= \mathbb P\bigl(p(\r)>\pcrit\bigr)
    \simeq \frac1{\sqrt{2\pi}\,z(\r)}\ee^{-z(\r)^2/2},
    \\
    z(\r) &= \frac{\pcrit}{\sigma_p(\r)}.
\end{align}
Evaluating this probability, and comparing it with the measured layering rate, is the
subject of \cref{sec:nucleation-results}.

\subsection{Layering probability and comparison with experiment}
\label{sec:nucleation-results}
% dean.tex 1392-1487 (remainder). Hero figure lives here.
The comparison with experiment is not to an absolute rate but to a \emph{ratio}:
Ref.~\cite{copenhagen2021} reports the distribution of distances between defects and new
layers, normalized by the same distribution for random points in the monolayer (their
Fig.~3f). The attempt rate is $p_\text{layer}/\tau_c$, and $\tau_c$ of \cref{eq:tau-c} is
essentially position-independent, so that ratio equals the ratio of layering
probabilities up to a slowly varying prefactor. \Cref{eq:p-layer} then gives
\begin{equation}
    \label{eq:R-ratio}
    \frac{p_\text{layer}(\r)}{p_\text{layer}(\infty)} \simeq
    \exp\!\left[\frac{z_\infty^2}2\left(1-\frac{\sigma_{p,\infty}^2}{\sigma_p^2(\r)}\right)\right],
    \qquad z_\infty \equiv \frac{\pcrit}{\sigma_{p,\infty}}.
\end{equation}
Only the \emph{relative} modulation of the variance matters, and the large prefactor
$z_\infty^2/2$ amplifies it. A modulation of a few percent therefore produces an
enhancement of two orders of magnitude: layer formation is exponentially sensitive to the
defect texture.

At the parameters of \cref{tab:used-values}, $\sigma_{p,\infty}=1.8\times10^3~
\Pa\cdot\micron$ and $z_\infty=19.1$. The variance modulation of \cref{fig:variance}
peaks at $\nu_0=3.1\%$, about $1.5~\micron$ off the core, and \cref{eq:R-ratio} then
gives a predicted enhancement
\begin{equation}
    \label{eq:R-value}
    \frac{p_\text{layer}(\r_\mathrm{max})}{p_\text{layer}(\infty)}
    = \exp\!\left[\frac{z_\infty^2}2\,\frac{\nu_0}{1+\nu_0}\right] \simeq 2.3\times10^2,
\end{equation}
against a first-bin value of $\simeq2\times10^2$ in Fig.~3f of Ref.~\cite{copenhagen2021},
whose reported error bar spans $40$--$500$ (\cref{fig:match}). The two agree within
that error bar.

The match is a one-parameter fit. Every quantity in \cref{eq:R-value} is fixed in advance
except the bulk modulus $B$, which \cref{sec:nucleation-energetics} already identified as
a model parameter rather than a measured constant. It enters the governing group
$\pcrit^2\elln^2\nu_0(\elln)/(B\zeta)$ only as $1/B$. \Cref{fig:match} should therefore be
read as a fit of $B$ and not as an independent prediction of the amplitude; the predicted
ratio is exponentially sensitive to the fixed inputs as well, most acutely to $\elln$ and
$h$ (\cref{app:extrusion-sensitivity}).

The \emph{shape} of the curve in \cref{fig:variance} is not fitted. The spatial profile
of $\sigmarho(\r)$ is the filtered double divergence of \cref{eq:total-gauss}, fixed once
$\Q_\ss$ is known, with no adjustable parameter: the decay over several micrometers, the
peak $1.5~\micron$ out on one side of the core, and the matching trough on the other. A
competing mechanism for the same observation would have to reproduce it. That shape, and
not the fitted amplitude of \cref{eq:R-value}, is the falsifiable content of the paper.

One normalization ambiguity in \cref{fig:match} remains open. Distance can be measured
from the point of maximum pressure variance, where layering is most probable, or from the
defect core, which is how the experimental histogram is binned but where the anisotropic
contribution to the variance is suppressed. The two curves shown bracket these readings.

\begin{figure}[tbp]
  \centering
  \includegraphics[width=\linewidth]{final.png}
  \caption{Predicted layering-probability enhancement, \cref{eq:R-ratio}, against the
    distance-binned ratio of Fig.~3f in Ref.~\cite{copenhagen2021} (error bar
    $40$--$500$). Two curves are shown: distance measured from the point of maximum
    pressure variance, and distance measured from the defect core itself, bracketing the
    two readings of the experimental binning convention. Computed with
    $B=2\times10^5~\Pa\cdot\micron$ and $\elln=0.5~\micron$, all other parameters as in
    \cref{tab:used-values}.}
  \label{fig:match}
\end{figure}

% ============================================================
\section{Discussion}
\label{sec:discussion}
% Pages ~1.5. New. structure.md \S4, four closing moves.
% ============================================================
The load-bearing part of this calculation is not the defect solution but the constraint
the coarse-graining imposes on it. Because the active stress and the noise both descend
from the same reversal term of \cref{eq:langevin-position}, they arrive with the same
coefficient $\zeta$, and the noise amplitude is not available to be tuned. Had the
equations been written down at the continuum level and a noise added to them, that
amplitude would have been a second adjustable quantity, and since the layering
probability depends on it exponentially, almost any observed enhancement could have been
reproduced. The comparison in \cref{sec:nucleation-results} carries weight only because
that freedom does not exist here.

The same rigidity makes the mechanism falsifiable. The threshold $\pcrit$ is fixed a
priori by adhesion, surface tension and cell geometry through \cref{eq:pcrit}, and
contains nothing that the comparison with experiment touches. Everything the reversal
rate controls sits in the other factor, through $\zeta\propto1/\frev$: slower reversals
mean longer runs, more active stress and louder noise together. Once $B$ has been fixed
at one reversal rate, the predicted enhancement at any other follows with nothing left to
adjust. Reversal-frequency mutants of the \textit{frz} pathway are standard tools, and
the sign is already on record: Ref.~\cite{han2025local} reports that non-reversing cells
show both larger traction fluctuations and more layer formation, which is what
$\zeta\propto1/\frev$ requires.
\todo{confirm the non-reversing (frzE) traction and layering observation is in
Ref.~\cite{han2025local} and not in a companion paper before this sentence stands.}

One parameter is fitted, the bulk modulus $B$, and it is worth being exact about why it
has to be. No measurement exists of the compressibility of a bacterial monolayer at the
scale that matters here, and the one apparent value in the literature for this system is
circular: the velocity fit of Ref.~\cite{copenhagen2021} constrains only the product of
$B$ with the cell extrusion flux, which is the quantity being computed in this paper. The
value we obtain, $B=2\times10^5~\Pa\cdot\micron$, lands close to the effective cell
modulus fitted to the monolayer-to-multilayer transition of \textit{E.\ coli}
\cite{grant2014}, a different organism undergoing the same kind of transition. Two fits
of the same kind of transition agreeing is a coincidence worth recording and not a
corroboration.

Several approximations limit how far the numbers should be pushed. The translational flux
of nematic order is dropped from both field equations, throughout the defect calculation
and not only in the steady state; the suppression that justifies it is
$(\ell_\star/\elld)^2\simeq2\times10^{-2}$ at the parameters used, but it is weakest at
the core, where the texture varies fastest, so the innermost part of the profile is the
least controlled. The analytic line of \cref{sec:defect-fluct} sets $\delta\Q=0$ and is
checked against the coupled numerical solution rather than against a bound. The
small-order-parameter closure limits the nematic order over which the implied noise
covariance stays positive semi-definite, and the defect steady state exceeds that limit
over much of its domain, which constrains the numerical solver rather than the analytic
result (\cref{app:realizability-failure,app:lyapunov}). The drift speed $u$ is not converged in
box size (\cref{app:comoving-results}). None of the adhesion or tension parameters of
\cref{sec:nucleation-energetics} has been measured for \textit{M.\ xanthus}, and the
predicted ratio is exponentially sensitive to several of them, most sharply to the layer
height $h$ (\cref{app:extrusion-sensitivity}). Finally, the monolayer is
treated as dry throughout, with no momentum conservation and therefore no backflow.

Two tests would go further than the one made here. A $-1/2$ defect generates no net
active flux and so neither drifts nor builds a density dipole, which makes its variance
signature qualitatively different rather than merely weaker; Ref.~\cite{copenhagen2021}
reports no layer formation there, and turning that null result into a quantitative
comparison is the natural next step. More directly, the variance field
\cref{eq:total-gauss} is itself an observable. Traction microscopy of the kind used in
Ref.~\cite{han2025local} measures the stress fluctuation, not only the layering events it
occasionally produces, so the prediction can be tested without passing through the
rare-event step at all, and without the fitted $B$, which cancels from the normalized
profile.

\begin{acknowledgments}
\skel{TODO.}
\end{acknowledgments}

% ============================================================
\appendix
\crefalias{section}{appendix}
\crefalias{subsection}{appendix}
% Appendices restructured 2026-09-22 to sit one-for-one under the main-text
% sections, in main-text order, so that every appendix has exactly one home:
%   A  -> \cref{sec:model-micro,sec:model-coarse,sec:model-closure}
%   B  -> \cref{sec:model-closure}   (the closure, and where it breaks)
%   C  -> parameters, first needed in \cref{sec:defect-steady}
%   D  -> \cref{sec:defect-steady}
%   E  -> \cref{sec:defect-fluct}
%   F  -> \cref{sec:nucleation}
% The realizability material, formerly a standalone appendix between B and D,
% is now B.5-B.7: it is a statement about the closure of appendix B and is
% never cited independently of it. The old cross-cutting "Parameters and
% supporting derivations" appendix is split, its model-parameter half becoming
% C and its extrusion half becoming F, so that \cref{sec:nucleation} gains the
% dedicated appendix it previously lacked.
% Every cross-reference goes through \cref, so appendix letters are never
% typed by hand.
% ============================================================

\section{Coarse-graining in full: from reversing rods to the field equations}
\label{app:dean}
% dean.tex 109-526.
This appendix gives in full the derivation compressed into
\cref{sec:model-micro,sec:model-coarse,sec:model-closure}, from the $N$-rod Langevin
system \cref{eq:model} to the stochastic field equations \cref{eq:model-rho,eq:model-Q}.

\subsection{Microscopic equations of motion}
\label{app:dean-microscopic}
\begin{itemize}
  \item Langevin equations for $N$ self-propelled rods: position, orientation,
    reversal (dichotomous noise $k_i$ at rate $\frev$, Gaussian-approximated on
    timescales $\gg \frev^{-1}$).
  \item The pairwise potential $V(\r-\r',\theta,\theta')$, its exchange and
    action-reaction symmetries, and the derivation of the force and torque fields
    \cref{eq:force-torque} from it.
  \item Full derivation of $\zeta = \xi_0 v_0^2/(4\frev)$ and of the force-density
    noise correlator, summarised without derivation in \cref{sec:model-micro}.
\end{itemize}

\subsection{Deriving the Dean equation for this model}
\label{app:dean-approach}
\begin{itemize}
  \item Full coarse-graining derivation, following \cite{dean1996} but for the
    orientation-resolved, self-propelled, reversing case: from the $N$-particle Langevin
    system to the field-theoretic Langevin equation for $\varrho_t(\r,\theta)$,
    multiplicative noise and all.
  \item The pair free energy \cref{eq:pair-free-energy} and the identification
    $\F=-\rho_0^{-1}\partial_\r\,\delta\calF/\delta\varrho$,
    $\Gamma=-\rho_0^{-1}\partial_\theta\,\delta\calF/\delta\varrho$.
  \item The noise sums, $\f=\rho_0^{-1}v_0\sum_i\varrho^{(i)}_t\p(\theta)k_i(t)$ and
    $g=\rho_0^{-1}\sum_i\varrho^{(i)}_t\lambda_i(t)$; the delta-function collapse that
    makes their correlators linear in $\varrho_t$; and the resulting
    $\sqrt{\varrho_t}$ multiplicative form, which is what lets the correlators be
    quoted as \cref{eq:noise-correlators} for $\varrho_t$-independent fields $k$,
    $\lambda$. State explicitly that matching second moments defines those fields only
    because the underlying noises have been taken Gaussian.
  \item State explicitly, since it is the crux, why the noise survives this
    coarse-graining where a naive continuum limit would drop it.
\end{itemize}

\subsection{Choosing the free energy functional}
\label{app:dean-free-energy}
\begin{itemize}
  \item Full justification of the incompressibility + Landau-de Gennes free energy
    density $\mathfrak f(\rho,\Q)$ (dean.tex eq:free-energy-choice-2), and of the
    one-constant approximation of the Frank elasticity.
  \item The chain rule \cref{eq:chain-rule} in full, including the functional
    derivatives $\delta\rho/\delta\varrho$ and $\delta\Q/\delta\varrho$ it rests on.
  \item State plainly what is \emph{not} done: $B$, $a$, $b$, $K$ are never expressed
    as moments of the microscopic potential $V$; $\mathfrak f$ is a modelling ansatz,
    not a computed coarse-grained limit of \cref{eq:pair-free-energy}.
\end{itemize}

\subsection{Normalisation of the nematic order parameter}
\label{app:dean-normalisation}
\begin{itemize}
  \item The two conventions: the conventional density-normalised
    $\Q=\rho^{-1}\!\int\!\d\theta\,\varrho\hat\Q$, bounded by unity, and the
    unnormalised moment \cref{eq:density-definition} used throughout this paper. They
    coincide at $\rho=1$.
  \item Why the unnormalised choice is made: with it, $\delta\Q/\delta\varrho =
    \hat\Q(\theta)\delta(\r-\r')$ exactly, so the chain rule \cref{eq:chain-rule}
    carries no extra $-\Q/\rho$ term and the angular integration closes cleanly.
  \item What has to be carried by hand as a consequence: the scalar order parameter is
    $S=|\Q|/\rho=\rho^{-1}\sqrt{(\Q\!:\!\Q)/2}$, the closure \cref{eq:closure} carries
    bare factors of $\rho$, and $\Q:\Q=2\rho^2S^2$. Give the reconciliation with the
    familiar $\Q=\rho S(2\n\n-\I)$ form.
  \item Note that the Pad\'e profile of \cref{sec:defect-steady} is written with the
    $\rho=1$ identification already made.
\end{itemize}

\subsection{Integrating out the angle}
\label{app:dean-angle}
\begin{itemize}
  \item Reduction from the joint density $\varrho_t(\r,\theta)$ to the density
    $\rho$ and nematic tensor $\Q$ fields, producing the stochastic active-nematic
    equations \cref{eq:model-rho,eq:model-Q} used from \cref{sec:model-closure} onward. The
    $\rho$ moment integrates exactly; the $\Q$ moment does not, and this is where the
    closure of \cref{app:closure} enters. State that dependency explicitly. State also
    why the polar moment $\int\!\d\theta\,\varrho_t\p(\theta)$ never enters.
  \item Explicit form of the density noise,
    $\bEta_\rho = \int\!\d\theta\,\sqrt{\varrho_t}\,\p(\theta)\,k(\r,\theta,t)$, and the
    contraction $\int\!\d\theta\,\varrho\,p_\alpha p_\beta=\tfrac12(Q_\ab+\rho\delta_\ab)$
    that turns \cref{eq:noise-correlators} into
    \cref{eq:eta-rho-autocorrelation}. Make explicit that $\bEta_\rho$ is a flux and not
    a force, so that it sits outside the mobility $1/\xi_0$ in \cref{eq:model-rho}, and
    that the rotational noise makes no contribution to the density channel.
  \item Explicit form and correlator of $\f^Q$, and its cross-correlator with
    $\bEta_\rho$; both descend from the same reversal noise $k$, so they are not
    independent. Give the isotropic rotational-noise part of $\f^Q$, of amplitude
    $8\Lambda\rho/\rho_0$, which is the one piece of the noise covariance that is local
    rather than a flux divergence.
  \item The large-$B$ truncation rule that selects which terms survive in
    \cref{eq:model-rho,eq:model-Q}: every term carrying a single factor $\nabla^n\rho$ is
    dropped unless $B$ multiplies it.
  \item Linearisation about a steady state, giving \cref{eq:linearized}, and the
    equivalent non-conservative form that separates the three channels by which the
    steady density reaches $\delta\Q$: it sources $\delta\Q$ from $\delta\rho$, advects
    $\delta\Q$, and amplifies it locally. Include the estimate showing why the frame
    term is negligible in the $\rho$ equation but retained in the $\Q$ equation.
\end{itemize}

\section{Closure of the angular moment hierarchy, and its limits}
\label{app:closure}
% dean.tex 2090-2262 (B.1-B.4), 2263-2450 (B.5-B.7).
This appendix supports the closure step of \cref{sec:model-closure}. The first four
subsections construct the maximum-entropy closure \cref{eq:closure}; the last three
establish where it stops being admissible, a limitation that constrains the numerical
solver of \cref{app:lyapunov} rather than the physics of \cref{eq:model-rho,eq:model-Q}.
The two halves were separate appendices in earlier drafts and are joined here because
the second is a statement about the first and is never used without it.

\subsection{The angular moments in integral form}
\label{app:closure-moments-integral}
\begin{itemize}
  \item Set-up of the angular-moment closure used to reduce the full orientation
    distribution to $\rho$, $\Q$ (feeds \cref{app:dean-angle}).
\end{itemize}

\subsection{Tensor basis}
\label{app:closure-tensor-basis}
\begin{itemize}
  \item The tensor basis ($\eye$, $\Q$, higher moments) the closure is expressed in.
\end{itemize}

\subsection{Exact moments}
\label{app:closure-exact-moments}
\begin{itemize}
  \item Exact expressions for the low-order angular moments before any small-parameter
    expansion.
\end{itemize}

\subsection{Small-$S$ expansion}
\label{app:closure-small-S}
\begin{itemize}
  \item Expansion of the exact moments at small nematic order $S$; recovers the
    truncated closure \cref{eq:closure} used in the main-line equations of
    \cref{sec:model-closure}, and states its $O(S^2)$ accuracy.
\end{itemize}

\subsection{The noise covariance is a Gram matrix}
\label{app:realizability-gram}
\begin{itemize}
  \item State the Gram-matrix structure and why it bounds what covariances are
    physically realizable.
\end{itemize}

\subsection{Where the closure ceases to be realizable}
\label{app:realizability-failure}
\begin{itemize}
  \item Where and how the closure of \cref{app:closure} fails at strong $S$: the noise
    covariance it implies is positive semi-definite only for
    $S\le S_\star=(\sqrt5-1)/2\simeq0.618$, and the defect steady state of
    \cref{sec:defect-steady} exceeds that bound over most of its domain.
  \item State the regime of failure in physical terms ($S$ range, not just ``strong''),
    so that the limitation quoted in \cref{sec:discussion} is checkable.
  \item State explicitly what is and is not affected: the $\delta\Q=0$ analytic result
    of \cref{sec:defect-fluct} is not, since its density-noise correlator
    \cref{eq:eta-rho-autocorrelation} does not depend on the closure; the full numerical
    solver of \cref{app:lyapunov} is.
\end{itemize}

\subsection{Restoring positivity with the exact moments}
\label{app:realizability-fix}
\begin{itemize}
  \item The positivity-preserving alternative using the exact (non-expanded) moments of
    \cref{app:closure-exact-moments}, at the cost of one scalar root-find per grid
    point; state whether this alternative was actually used anywhere in the main line
    or is offered as a fallback. It currently reads as the latter; confirm.
\end{itemize}

\section{Model parameters}
\label{app:parameters}
% dean.tex 543-680 (C.1), 2451-2459 (C.2, folded in).
% Formerly the first two thirds of a cross-cutting "Parameters and supporting
% derivations" appendix; its extrusion third is now \cref{app:extrusion}.
The parameters of \cref{eq:model-rho,eq:model-Q} are first needed in
\cref{sec:defect-steady} and are used unchanged from there on. The separate set of
adhesion, tension and geometry parameters that enters the nucleation argument of
\cref{sec:nucleation} is collected in \cref{app:extrusion}.

\subsection{Parameter choice}
\label{app:parameters-choice}
\begin{itemize}
  \item Full model-parameter table: ``known'' group from experiment
    (\cite{copenhagen2021,han2025local} and project estimates), ``guessed'' group
    ($B$, nematic relaxation $\gamma$, $\Lambda$), derived group ($\xir$, bare $a,b,K$).
  \item State the bare-vs.-activity-renormalised distinction carefully: $a'=a+\Lambda\xir$
    and $K'=K+\zeta\xir/(4\xi_0)$ are what set $\elld$ and $S_0$, not the bare $a$, $K$.
  \item \todo{this table has moved at least three times during the project (see
    project memory: canonical-model-parameters); pull the values fresh from dean.tex
    at time of writing, do not copy from an earlier draft or a jobscript.}
\end{itemize}

\subsection{Relaxation rate of the nematic order parameter}
\label{app:parameters-relaxation}
\begin{itemize}
  \item Short, self-contained derivation of $\gamma = 8b/\xir$ (dean.tex
    app-Q-relaxation). It feeds \cref{app:parameters-choice} only, and is kept as a
    subsection here rather than a standalone appendix since nothing else cites it.
\end{itemize}

\section{Comoving-frame defect solver}
\label{app:comoving}
% dean.tex 2460-2597.
This appendix supports \cref{sec:defect-steady}: it is the numerical route to the
deterministic steady state $\rho_\ss$, $\Q_\ss$ and to the drift speed $u$, together
with the closed-form companion solution quoted there.

\subsection{The comoving-frame equations and their discretisation}
\label{app:comoving-disc}
Dropping the $\nabla\H$ terms as in \cref{sec:model-closure}, regrouping the remaining
local terms of \cref{eq:model-Q} into $\Ha$ of \cref{eq:H-active}, and substituting
$\partial_t\to\partial_t-(\u\cdot\nabla)$, the noiseless steady state solves
\begin{subequations}
    \label{eq:comoving-equations}
\begin{align}
    0 ={}& \u\!\cdot\!\nabla\rho + \frac1{\xi_0}\nabla\!\cdot\!
    \bigl(B\nabla\rho - \fa\bigr),
    \\
    0 ={}& \nabla\!\cdot\!\Bigl[\Q\Bigl(\u+\frac B{\xi_0}\nabla\rho\Bigr)\Bigr]
    + \frac4{\xir}\Ha .
\end{align}
\end{subequations}
The frame velocity $\u$ is not available in closed form. It is fixed implicitly, by
requiring that the defect core stay pinned at the origin, and enters the discretized
problem as the Lagrange multiplier conjugate to that constraint. Both fields are held at
their far-field values on the boundary of the computational domain, a Dirichlet condition
on $\Q$ and on $\rho$ alike, which corresponds to a window cut from a larger colony. The
sealed-wall alternative and what changes under it are given below and in
\cref{app:lyapunov-bc}.

\Cref{eq:comoving-equations} is solved by finite differences on a graded, non-uniform
grid, fourth order in the local spacing and refined toward the defect core.
\begin{itemize}
  \item Graded non-uniform grid, fourth-order finite differences, Kronecker-lifted to
    2D (full grid-construction algorithm here).
  \item Boundary conditions in full. $\Q$ is held at the Pad\'e $+1/2$ profile on the
    whole boundary in every run. For $\rho$ there are two cases: the reservoir
    (Dirichlet) wall $\rho=1$ used in the main text, and the sealed wall
    $\nh\cdot(B\nabla\rho+\zeta\nabla\!\cdot\!\Q)=0$. Note that the sealed condition is
    on the \emph{total} flux, not on $\nabla\rho$ alone, that the normal must be signed
    at the corners, and that the sealed problem needs one extra condition to fix the
    constant mode. \Cref{fig:rho1d-dirichlet-neumann} compares the two solutions.
  \item $\u$ as a Lagrange multiplier conjugate to the pinning constraint: full
    bordered-system derivation, and why the closed-form core condition
    $\u\simeq-(B/\xi_0)\nabla\rho(0)$ is a poor way to \emph{compute} $\u$, overestimating
    $|\u|$ by a factor $3.75$ at the reservoir wall.
  \item Practical notes worth keeping (each cost real development time per the source
    material): row equilibration at large $B$, continuation in $\zeta$ from zero,
    constraint-row scaling for the differential-algebraic system.
  \item The closed-form Pad\'e-$\Q$ route to $\rho_\ss$ used as the analytic companion
    in \cref{sec:defect-steady}: the auxiliary radial functions, and the coefficients
    under both Dirichlet and Neumann far-field conditions on a disc.
\end{itemize}

\subsection{Verification}
\label{app:comoving-verification}
\begin{itemize}
  \item Full verification table: analytic-vs.-numerical Jacobian, passive-limit checks
    against an independent 1D BVP solve, exact pinning, mass/multiplier conservation,
    mirror-symmetry residual, lab-frame-vs.-bordered-$\u$ cross-check.
  \item The passive-limit row, $u\to0$ as $\zeta\to0$, is the numerical statement that
    the drift of \cref{sec:defect-steady} is activity-driven; quote its convergence
    order.
  \item State which single row is the decisive one (lab-frame drift vs.\ bordered $\u$,
    per the source material) and why: it uses none of the border/constraint machinery,
    so agreement there is two independent numerical routes to the same number.
\end{itemize}

\subsection{Results}
\label{app:comoving-results}
\begin{itemize}
  \item Full results beyond the summary in \cref{sec:defect-steady}: transient
    overshoot in $u(t)$, its physical origin (density dipole not yet built at $t=0$).
  \item Both wall velocities, and the accuracy of the closed-form core estimate against
    each (a factor $3.75$ at the reservoir wall, $2.12$ at the sealed one).
  \item Box-size convergence, moved here from the main text: $u$ is \emph{not}
    converged in box size, changing by $10.8\%$ between $L=8$ and $16~\micron$, and the
    reservoir and sealed-wall conditions bracket the true value from opposite sides
    rather than pinning it, with a gap closing like $1/L$. The grid error is negligible
    by comparison, so it is the box and not the mesh that limits the quoted value.
  \item Sealed-vs.-reservoir comparison in full, including the unresolved cost
    difference between the two (sealed wall $\sim7\times$ more expensive per solve,
    reason not currently understood). State this as an open numerical question, not a
    physics one.
\end{itemize}

\begin{figure}[tbp]
  \centering
  \includegraphics[width=0.85\linewidth]{strength.png}
  \caption{Nematic order $S$ against distance from the core in the converged
    comoving-frame solution, along $+x$, along $-x$ and along $\pm y$, compared with the
    Pad\'e approximation \cref{eq:pade-approximation} (dashed). The coupling to the
    density splits the $\pm x$ curves, which is the distortion the drift produces; the
    departure from the approximant stays small, which is what licenses using it in
    \cref{sec:defect-steady}.}
  \label{fig:strength}
\end{figure}

\begin{figure}[tbp]
  \centering
  \includegraphics[width=0.85\linewidth]{defectangle.png}
  \caption{Director angle in the converged solution minus the ideal $+1/2$ profile,
    $\theta-\phi/2$, against polar angle $\phi$ at four radii. The deviation changes sign
    between the two sides of the defect axis, grows with radius, and reaches one degree
    at $r=5~\micron$.}
  \label{fig:defectangle}
\end{figure}

\begin{figure}[tbp]
  \centering
  \includegraphics[width=0.85\linewidth]{activeforce.png}
  \caption{Active force along the defect axis in the converged solution, compared with
    \cref{eq:active-force-defect} evaluated on the Pad\'e profile. The small deviation
    splits the two sides of the core, as in \cref{fig:strength}, and has the same origin
    in the drift.}
  \label{fig:activeforce}
\end{figure}

\begin{figure}[tbp]
  \centering
  \includegraphics[width=0.85\linewidth]{rho1d_Dirichlet_Neumann.png}
  \caption{Boundary-condition dependence of the steady-state density, same $y=0$ section
    as \cref{fig:rho1d-dirichlet}. Blue: reservoir (Dirichlet) wall, $\rho=1$ on the
    boundary, used throughout the main text. Orange: sealed (Neumann) wall, zero total
    flux. Dashed curves are the corresponding closed-form solutions. The two agree near
    the core and differ in the far field, where the sealed wall cannot relax the
    accumulated density, which is also why the two walls give different drift speeds
    (\cref{app:comoving-results}).}
  \label{fig:rho1d-dirichlet-neumann}
\end{figure}



\section{Lyapunov solver for the coupled fluctuations}
\label{app:lyapunov}
% dean.tex 1521-2089.
This appendix supports \cref{sec:defect-fluct}: it is the numerical route to the
equal-time covariance of $(\delta\rho,\delta\Q)$ on the steady background of
\cref{app:comoving}, retaining the coupling to $\delta\Q$ that the analytic result
\cref{eq:total-gauss} drops, and it derives the short-length regularisation
\cref{eq:cutoff} that both routes share.

\subsection{Steady-state covariance via the Lyapunov equation}
\label{app:lyapunov-eq}
\begin{itemize}
  \item General linear OU field theory set-up: $\partial_t\psi = \A(\r)\psi+\bxi$,
    steady-state covariance solves the Lyapunov equation. State this generally before
    specialising to $\psi=(\delta\rho,\delta\Q)$.
\end{itemize}

\subsection{Discretisation}
\label{app:lyapunov-disc}
\begin{itemize}
  \item Finite-difference discretisation of $\A$ and the noise covariance $\B$; how the
    Lyapunov equation is solved numerically for the discretised system.
  \item Where the closure of \cref{app:closure} enters the noise covariance, and the
    practical consequence of the realizability bound of
    \cref{app:realizability-failure}: the solver is asked for a covariance the closure
    does not guarantee to be positive semi-definite over most of the defect domain.
    Report what is done about it, and what the exact-moment alternative of
    \cref{app:realizability-fix} would cost.
\end{itemize}

\subsection{Why the choice of difference operators matters}
\label{app:lyapunov-operators}
\begin{itemize}
  \item Self-adjointness of the discretised operators vs.\ the continuum ones, and the
    numerical artefacts that appear if this is gotten wrong. Keep this appendix even
    though it is ``just numerics,'' since it is what makes the boundary results of
    \cref{app:lyapunov-bc} trustworthy.
\end{itemize}

\subsection{Boundary conditions}
\label{app:lyapunov-bc}
\begin{itemize}
  \item Reservoir vs.\ sealed-wall boundary conditions for the fluctuation solver, and
    what changes in the covariance near the boundary.
  \item Spell out that this is not one condition but three: on the steady state, on the
    fluctuating field, and on the noise flux through the wall. Give the table of the
    eight combinations and show that only two are physically coherent, the reservoir
    (all Dirichlet, noise flux free) that the solver runs and used in
    \cref{sec:defect-steady}, and the sealed one (all Neumann, noise flux blocked).
  \item Show why the isotropic pedestal of \cref{eq:total-gauss} is blind to the choice,
    and why the anisotropic wall residue is not: it reverses sign and grows by roughly
    a factor of three between the two, a no-flux wall being quiet where a reservoir
    wall is loud.
  \item Report the active-force fluctuations the solver returns,
    \cref{fig:fxfx,fig:fyfy}, and what they check.
\end{itemize}


\subsection{The short-length cutoff and noise with a finite correlation length}
\label{app:lyapunov-finite-corr}
\begin{itemize}
  \item How the cutoff $\elln$ of \cref{eq:cutoff} is implemented: the fields are
    smoothed over a Gaussian patch, so the two-point function carries the kernel twice
    and the effective smoothing of the correlator is a Gaussian of width $\elln$.
  \item Derivation of the resulting variance \cref{eq:total-gauss}: the isotropic
    pedestal, and the anisotropic term as a filtered double divergence with Fourier
    filter $\hat\Phi_{\elln}(p)=\tfrac12E_1(p^2\elln^2/8)$, its two asymptotic branches,
    and its real-space transform \cref{eq:kappa-real}.
  \item State that the filter involves no expansion and no fitted scale: the part of the
    internal-wavevector integral quadratic in that wavevector vanishes by isotropy
    against the traceless $\Q_\ss$, and the exponential the remainder produces is
    cancelled exactly by the external smoothing factor.
  \item The independent hard-cutoff (Duhamel) derivation that reaches the same
    pedestal exactly and the same leading logarithm, as a check on the Gaussian
    regularisation.
  \item Extension of the white-noise treatment to a noise kernel with a genuinely
    finite correlation length; spell out whether that is the same regularisation as the
    resolution cutoff above or a different one. This was not fully resolved as of this
    outline.
\end{itemize}

\begin{figure}[tbp]
  \centering
  \includegraphics[width=0.85\linewidth]{ln_curve.png}
  \caption{Peak variance modulation $\nu_0$ against the cutoff $\elln$. The modulation
    vanishes as $\elln\to0$: without a short-length cutoff the defect leaves no signature
    in the density variance. The operating point is $\elln=0.5~\micron$
    (\cref{eq:cutoff}), the width of an \textit{M.\ xanthus} cell.}
  \label{fig:ln_curve}
\end{figure}

\begin{figure}[tbp]
  \centering
  \includegraphics[width=0.85\linewidth]{fxfx.png}
  \caption{%
\capitem Active-force fluctuation, $x$-component variance, from the Lyapunov solver:
      supporting evidence for the validation claim of \cref{sec:defect-fluct}.
}
  \label{fig:fxfx}
\end{figure}

\begin{figure}[tbp]
  \centering
  \includegraphics[width=0.85\linewidth]{fyfy.png}
  \caption{%
\capitem Companion $y$-component, same purpose as \cref{fig:fxfx}.
}
  \label{fig:fyfy}
\end{figure}

\section{Extrusion energetics: parameters, ranges and sensitivity}
\label{app:extrusion}
% dean.tex 1329-1360, plus the sensitivity material moved out of
% \cref{sec:nucleation-results}. Split out of the former cross-cutting
% "Parameters and supporting derivations" appendix so that
% \cref{sec:nucleation} has an appendix of its own; the model parameters
% it left behind are \cref{app:parameters}.
This appendix supports \cref{sec:nucleation}. The main text quotes one value per
parameter in \cref{tab:used-values} and one sensitivity claim; everything behind both
is here.

\subsection{Literature ranges for the extrusion parameters}
\label{app:parameters-extrusion}
\begin{itemize}
  \item Full literature table for $\gammaw$, $\Wc$, $\Ws$, $B$, $\elln$, $h$, $\Acell$:
    plausible ranges and origin for each, as currently in dean.tex \S``Parameter
    estimates and their ranges''. This is what \cref{sec:nucleation-energetics}
    compresses to ``used values only.''
  \item The status of each entry. $B$ is the one fitted quantity. Every other entry,
    $\elln$ included, is a model parameter fixed before the comparison with experiment.
    Record that the rod second-moment estimate
    $\ell_\mathrm{rod}/\sqrt{12}\simeq\sqrt3~\micron$ is an upper bound only, that it
    would violate the scale-separation requirement $\elln\ll\elld$, and that $\elln$ and
    the layer height $h$ are numerically the same length but enter through different
    physics.
\end{itemize}

\subsection{How the prediction depends on the parameters}
\label{app:extrusion-sensitivity}
\begin{itemize}
  \item Full one-at-a-time sensitivity analysis for the predicted enhancement ratio,
    moved here from \cref{sec:nucleation-results}: each parameter varied about its
    used value, with multipliers currently spanning $\times0.013$ to
    $\times2.3\times10^{2}$. Pull current numbers from dean.tex rather than hard-coding
    them here; they are moving.
  \item The $\elln$ dependence specifically, since it is the steepest: with
    $z_\infty\propto\elln$ and $\nu_0\propto\elln^2$ at small $\elln$, the exponent grows
    as a high power of $\elln$, the effective power near the operating point being close
    to $\elln^{3.3}$. Halving $\elln$ takes the predicted ratio from $2.3\times10^2$ to
    below $2$; raising it to $0.7~\micron$ takes it past $10^6$. Cross-reference
    \cref{fig:ln_curve}, which shows $\nu_0(\elln)$ vanishing as $\elln\to0$.
\end{itemize}

\subsection{The bulk modulus as the single fitted parameter}
\label{app:extrusion-B}
\begin{itemize}
  \item State the fit argument in full here (the compressed version lives in
    \cref{sec:nucleation-results}): the single governing group
    $\ln(\text{ratio}) = \tfrac{\pcrit^2}{2}\cdot\tfrac{2\pi\rho_0\elln^2}{B\zeta}\cdot
    \tfrac{\nuzero}{1+\nuzero}$, and why $B$ enters only as $1/B$.
  \item Note the $B$ ladder explicitly, including the coincidence flagged in
    \cref{sec:discussion}: colony rheology \cite{kochanowski2024} at the low end, the
    \textit{E.\ coli} mono-to-multilayer modulus \cite{grant2014} near where the fitted
    $B$ lands, cell-envelope Young's modulus as an implausible ceiling. Note also that
    the apparent $B$ obtainable from Ref.~\cite{copenhagen2021} is circular, since
    their velocity fit constrains only the product of $B$ with the extrusion flux this
    paper is computing.
  \item Cross-reference \cref{app:parameters-choice}, where the same $B$ enters the
    field equations, so that the reader meets one fitted number and not two.
\end{itemize}

\subsection{The energetics idealizations in full}
\label{app:extrusion-idealizations}
\begin{itemize}
  \item The two bounding idealizations of \cref{sec:nucleation-energetics} in full: the
    $A$-vs.-$2A$ substrate-detachment bracket and its renormalisation of $\Ws$ and $B$,
    and the partial-wetting bound on $\gammaw$.
  \item The subleading terms of \cref{eq:dEe-full} that \cref{eq:dEe} drops, and the
    size of the correction at the operating compression $\epsilon=\pcrit/B$.
\end{itemize}

\bibliographystyle{apsrev4-2}
\bibliography{literature}

\end{document}
